%%% Set margins, indents, and linespread %%%

\usepackage[left=30mm, right=15mm, top=20mm, bottom=20mm]{geometry} % Package arguments define page margins
\pagestyle{plain} % Remove default upper header and leave only a page number at the bottom center of the page
\parindent=1.25cm % First line indent 1.25 cm, approximately equal to 5 characters (as required by GOST)
\usepackage{indentfirst} % Add indent to the first paragraph
%\linespread{1.3} % Linespread (the closest to 1.5 in Word) - here we use a command OnehalfSpacing* later


%%% Set language parameters and font %%%

\usepackage[russian]{babel}                         % English is the only language for the document
\usepackage{siunitx}
\babelfont{rm}{Times New Roman}    % TMR as the primary font
\babelfont{tt}{Courier New}    % cyrilic monofont 
\sisetup{output-decimal-marker = {,}}
\usepackage{fontenc}
\usepackage{fontspec} %ru in verbatim
\usepackage{verbatim}
%%% Set heading style for chapters, sections, and subsections %%%

% \usepackage{calc} % for indentation in the table of contents (before "c.")

\setsecnumdepth{subsection} % Numbers are only assigned to the first three levels, i.e., numbered are chapters, sections, and subsections
\renewcommand*{\chapterheadstart}{} % Redefine the command which adds spacing above the heading to prevent it from doing this
\renewcommand*{\printchaptername}{\ifbool{cent}{}{\hspace{1.25cm}}} % Redefine the command which prints the word 'Chapter' to prevent it from doing this
%\renewcommand*{\printchapternum}{} % Same for the chapter number - no, here we need to keep the number
\renewcommand*{\chapnumfont}{\normalfont\bfseries} % Change the chapter number font style: normal size, bold
\renewcommand*{\afterchapternum}{\hspace{1em}} % Change the delimiter between the chapter number and the chapter title
\renewcommand*{\printchaptertitle}{\normalfont\bfseries\ifbool{cent}{\centering}{}} % Change the chapter title font style: normal size, bold, centered, uppercase

\renewcommand*{\printtoctitle}[1]{%
  \begin{center}
    \normalfont\bfseries СОДЕРЖАНИЕ
  \end{center}
  \vspace{-1.3pt} \vspace{-\cftbeforechapterskip}
  \hspace{\linewidth} \hspace{-19pt} \textbf{с.}
} % unbelievable hack for indentation in the table of contents (before "c.")

\setlength{\beforechapskip}{1.3pt}
\setlength{\afterchapskip}{1.3pt}


\setbeforesecskip{1.3pt} % Set spacing above the section heading 
\setaftersecskip{1.3pt} % Same spacing below the section heading
\setsecheadstyle{\raggedright\hspace{1.25cm}\normalfont\bfseries} % Change the section heading font style: right aligned, normal size, bold
\setbeforesubsecskip{1.3pt} % Set spacing above the subsection heading
\setaftersubsecskip{1.3pt} % Same spacing below the subsection heading
\setsubsecheadstyle{\raggedright\normalfont\bfseries}  % Change the subsection heading font style: right aligned, normal size, bold



\setrmarg{2.55em plus1fil} % Disable hyphenation in toc
\setlength{\cftbeforechapterskip}{1.3pt} % This command removes vertical spacing before chapter titles - no, keep the spacing
% \renewcommand{\aftertoctitle}{\afterchaptertitle \vspace{-1.3pt} \vspace{-\cftbeforechapterskip} \begin{flushleft}\hfill с.\end{flushleft} 
% \vspace{-3.9pt}
% \vspace{-\cftbeforechapterskip}
% } % Adjust vertical spacing right after the heading 'Table of Contents' to make it the same as for other chapter headings
% \renewcommand{\aftertoctitle}{\\[0.5\baselineskip]\hfill\textnormal{с.}}
\renewcommand*{\chapternumberline}[1]{} % Redefine the command which prints the chapter number to prevent it from doing this - no, keep the number here
% \renewcommand*{\cftchapternumwidth}{1.5em} % Change the delimiter between the chapter number and the chapter title
\renewcommand*{\cftchapterfont}{\normalfont} % Chapter titles are printed with normal size, uppercase
\renewcommand*{\cftchapterpagefont}{\normalfont} % Page numbers are printed with normal size
\renewcommand*{\cftchapterdotsep}{\cftdotsep} % Create dot separators which go up to the page number for each entry in toc
\renewcommand*{\cftdotsep}{0} % Define the dot frequency
\renewcommand*{\cftchapterleader}{\cftdotfill{\cftchapterdotsep}} % Make the dots of standard shape (by default they are 'bold')
\maxtocdepth{section} % Only the first three levels are displayed in toc: chapters, sections, and subsections



%%% Alignment and hyphenation %%%

%% http://tex.stackexchange.com/questions/241343/what-is-the-meaning-of-fussy-sloppy-emergencystretch-tolerance-hbadness
%% http://www.latex-community.org/forum/viewtopic.php?p=70342#p70342
\tolerance 1414
\hbadness 1414
\emergencystretch 1.5em                             % If something goes wrong, try to adjust this one first
\hfuzz 0.3pt
\vfuzz \hfuzz
%\dbottom
%\sloppy                                            % No overfilling
\clubpenalty=10000                                  % Prevent page break after the first line in a paragraph
\widowpenalty=10000                                 % Prevent page break before the last line in a paragraph
\brokenpenalty=4991                                 % Prevent page break directly after hyphenation

%%% Set parameters for figures and tables %%%

\usepackage{graphicx, caption, subcaption} % Load packages for figures and captions
\graphicspath{{images/}} % Define the image folder
\captionsetup[figure]{font=normal, width=\textwidth, name=Рисунок, justification=centering, parskip=6pt} % Set the figure caption parameters: small font size (12pt), caption width is equal to the width of the main text, complete word "Figure", caption is centered
\captionsetup[subfigure]{font=normal} % Subfigure indices a), b), ... are also in small size (12pt) (by default they are even smaller)
\captionsetup[table]{singlelinecheck=false,font=normal,width=\textwidth,justification=justified, parskip=4pt, name=Таблица} % Set the table caption parameters: disable hyphenation, small font size (12pt), caption width is equal to the width of the main text, caption is justified
\captiondelim{ --- } % Set em dash as the delimiter between the caption name (i.e., the words 'Figure' and 'Table') and caption text
\setkeys{Gin}{width=\textwidth} % By defalut, the size of all figures will be adjusted to fit into the main text
%\setlength{\abovecaptionskip}{0pt} % Vertical space above the caption - no, keep the default
%\setlength{\belowcaptionskip}{0pt} % Vertical space below the caption - no, keep the default
\usepackage[section]{placeins} % All objects of the float type (figures/tables) do not go beyond the section where they are initialized

%%% Set parameters for references and hyperreferences %%%

\usepackage[unicode=ture]{hyperref}                               % Load the corresponding package
\hypersetup{
    colorlinks=true,                                % All refs and hyperrefs appear colored in text
    linktoc=all,                                    % Clickable colored refs for all entries in toc
    linktocpage=false,                               % Only the page numbers will be colored in toc, not the whole titles
    linkcolor=black,                                  % The color of refs and hyperrefs is red
    citecolor=black,                                 % The color of citations is red
    urlcolor = black
}

\usepackage{bookmark}
\bookmarksetup{
  numbered=true,    
  open=true,        
  depth=2            
}

%%% Set list parameters %%%

\usepackage{enumitem}                               % Load the package for the fine adjustment of list appearance
\renewcommand*{\labelitemi}{\normalfont{--}}        % Use en dash for unnumbered list items
\setlist{noitemsep, leftmargin=*,
%nolistsep
%,itemindent=\parindent
}                   % Remove vertical spacings between the items of the same level
\setlist[1]{labelindent=\parindent}                 % Item indent is equal to the paragraph indent
%\setlist[2]{leftmargin=\parindent}                  % Additional indent for the 2nd level
%\setlist[3]{leftmargin=\parindent}                  % And one more for the 3rd level

%%% Counters and object numbering %%%

\counterwithout{figure}{chapter}                    % Continuous numbering of figures throughout the document
\counterwithout{equation}{chapter}                  % Continuous numbering of math expressions throughout the document
\counterwithout{table}{chapter}                     % Continuous numbering of tables throughout the document

%%% Create bibliography using the biblatex package and biblatex-gost style processed by biber %%%

\usepackage{csquotes} % This package manages complex citation blocks (biblatex recommends to load it)
\usepackage[%
backend=biber,  % Processor
bibencoding=utf8,                                   % Encoding of the bib-file
sorting=none,                                       % Preference for sorting the reference list
style=gost-numeric,                                 % Citation style (GOST)
language=russian,                                   % Russian language for the reference list
sortcites=true,                                     % Several cites in square brackets will appear sorted
movenames=false,                                    % Do not move author names, they are always in the beginning of the citation entry
maxnames=5,                                         % Maximum number of authors for an entry
minnames=3,                                         % If there is more that 'max', cut the number of authors up to this
doi=false,                                          % Do now show DOI
isbn=false,    % Do not show ISBN, ISSN, ISRN
locallabelwidth=true,

]{biblatex}
\DeclareDelimFormat{bibinitdelim}{}                 % Remove the space between initials (Readdler J.R. instead of Readdler J. R.)
\addbibresource{bibl.bib}                           % Define the bibliography file



%%% Additional packages for enhanced functionality %%%

\usepackage{longtable,ltcaption}                    % Long tables
\usepackage{multirow,makecell}                      % Enriched formatting for tables
\usepackage{booktabs}                               % Yet another package for nice tables
\usepackage{soulutf8}                               % Support hyphenatable letterspacing, underline, and strikethrough
\usepackage[none]{hyphenat}                               % Enhanced hyphenation functionality
\usepackage{textcomp}                               % Support for 'complex' printed signs like some currency symbols, copyright, etc.
\usepackage[version=4]{mhchem}                      % Beautiful chemical equations
\usepackage{amsmath} % Enriched appearance of mathematical expressions 
\usepackage{amsthm}
\numberwithin{equation}{chapter}  % нумерация ур-ний вида 1.1, 1.2 (если chapter = 1)

%%% Subsequently add all document parts %%%

\setlength{\parskip}{0pt}


\usepackage{pdfpages}
\usepackage{float}
\usepackage{amssymb} % Для \mathbb{R}
\usepackage{tikz}
\usetikzlibrary{arrows.meta}
\usepackage[e]{esvect} %nice vec
\usepackage{mathtools} %rcases

\newtheorem{theorem}{Теорема}
\renewcommand\qedsymbol{$\blacksquare$}

\newcommand{\R}{\mathbb{R}}


\newbool{cent}